% GENERATED FILE. Edit canonical lessons, then run npm run generate:books.
% canonical-source-hash: fnv1a64:c4d7c79791b31fbb
% canonical-lessons: ML-C144-atukontu, ML-C144-uddesyam, ML-C144-phalam, ML-C144-hetu, ML-C144-nimittam

\chapter{So, Purpose, Result, Cause, On account of}
\label{ch:ml-so-purpose-result-cause-on-account-of}
\hlchaptermodality{144}

\begin{chapteropening}
\textbf{By the end of this chapter:} \emph{I can say so, a purpose, a result, a cause and on account of.}

Complete the last lesson of chapter 144: i can say so, a purpose, a result, a cause and on account of.
\end{chapteropening}

\section[\texorpdfstring{\ml{അതുകൊണ്ട്} (atukoṇṭŭ)}{അതുകൊണ്ട് (atukoṇṭŭ)}]{\ml{അതുകൊണ്ട്} (atukoṇṭŭ) — so, therefore}

\label{lesson:ML-C144-atukontu}



\begin{quote}
Before the new one: say the Malayalam for a habit, then the Malayalam for a qualification.
\end{quote}

\subsection*{You'll want to know: \ml{അതുകൊണ്ട്}}
\textbf{\ml{അതുകൊണ്ട്}} — \emph{atukoṇṭŭ} — \textquotedblleft{}so, therefore\textquotedblright{}.

\begin{grammarlens}[title={What you've built}]
One more word for this chapter.
\end{grammarlens}

\subsection*{Your turn}
\begin{itemize}
\raggedright
  \item \emph{Say it:} \emph{atukoṇṭŭ}
  \item \emph{Say it:} \emph{atukoṇṭŭ}, once more
  \item \emph{Say it:} say \emph{atukoṇṭŭ}
  \item \emph{Recall:} say the Malayalam for a habit, then the Malayalam for a qualification, then say \emph{atukoṇṭŭ} again
\end{itemize}

\begin{tcolorbox}[breakable,colback=teal!4,colframe=teal!35!black,title={Before you move on}]
What does \ml{അതുകൊണ്ട്} mean? (\textquotedblleft{}So, therefore\textquotedblright{}.) Say it once more.
\end{tcolorbox}
\section[\texorpdfstring{\ml{ഉദ്ദേശ്യം} (uddēśyaṁ)}{ഉദ്ദേശ്യം (uddēśyaṁ)}]{\ml{ഉദ്ദേശ്യം} (uddēśyaṁ) — a purpose, an intention}

\label{lesson:ML-C144-uddesyam}



\begin{quote}
Before the new one: say the Malayalam for a qualification, then the Malayalam for so.

Two older, literary words for night: \emph{iravŭ} and \emph{iruḷ}. Which one is also the Tamil word? (\emph{Iravŭ}.)
\end{quote}

\subsection*{You'll want to know: \ml{ഉദ്ദേശ്യം}}
\textbf{\ml{ഉദ്ദേശ്യം}} — \emph{uddēśyaṁ} — \textquotedblleft{}a purpose, an intention\textquotedblright{}.

\begin{grammarlens}[title={What you've built}]
One more word for this chapter.
\end{grammarlens}

\subsection*{Your turn}
\begin{itemize}
\raggedright
  \item \emph{Say it:} \emph{uddēśyaṁ}
  \item \emph{Say it:} \emph{uddēśyaṁ}, once more
  \item \emph{Say it:} say \emph{uddēśyaṁ}
  \item \emph{Recall:} say the Malayalam for a qualification, then the Malayalam for so, then say \emph{uddēśyaṁ} again
\end{itemize}

\begin{tcolorbox}[breakable,colback=teal!4,colframe=teal!35!black,title={Before you move on}]
What does \ml{ഉദ്ദേശ്യം} mean? (\textquotedblleft{}A purpose, an intention\textquotedblright{}.) Say it once more.
\end{tcolorbox}
\section[\texorpdfstring{\ml{ഫലം} (phalaṁ)}{ഫലം (phalaṁ)}]{\ml{ഫലം} (phalaṁ) — a result}

\label{lesson:ML-C144-phalam}



\begin{quote}
Before the new one: say the Malayalam for so, then the Malayalam for a purpose.
\end{quote}

\subsection*{You'll want to know: \ml{ഫലം}}
\textbf{\ml{ഫലം}} — \emph{phalaṁ} — \textquotedblleft{}a result\textquotedblright{}.

\begin{grammarlens}[title={What you've built}]
One more word for this chapter.
\end{grammarlens}

\subsection*{Your turn}
\begin{itemize}
\raggedright
  \item \emph{Say it:} \emph{phalaṁ}
  \item \emph{Say it:} \emph{phalaṁ}, once more
  \item \emph{Say it:} say \emph{phalaṁ}
  \item \emph{Recall:} say the Malayalam for so, then the Malayalam for a purpose, then say \emph{phalaṁ} again
\end{itemize}

\begin{tcolorbox}[breakable,colback=teal!4,colframe=teal!35!black,title={Before you move on}]
What does \ml{ഫലം} mean? (\textquotedblleft{}A result\textquotedblright{}.) Say it once more.
\end{tcolorbox}
\section[\texorpdfstring{\ml{ഹേതു} (hētu)}{ഹേതു (hētu)}]{\ml{ഹേതു} (hētu) — a cause}

\label{lesson:ML-C144-hetu}



\begin{quote}
Before the new one: say the Malayalam for a purpose, then the Malayalam for a result.

Say the two literary night words once more. (\emph{Iravŭ}, \emph{iruḷ}.)
\end{quote}

\subsection*{You'll want to know: \ml{ഹേതു}}
\textbf{\ml{ഹേതു}} — \emph{hētu} — \textquotedblleft{}a cause\textquotedblright{}.

\begin{grammarlens}[title={What you've built}]
One more word for this chapter.
\end{grammarlens}

\subsection*{Your turn}
\begin{itemize}
\raggedright
  \item \emph{Say it:} \emph{hētu}
  \item \emph{Say it:} \emph{hētu}, once more
  \item \emph{Say it:} say \emph{hētu}
  \item \emph{Recall:} say the Malayalam for a purpose, then the Malayalam for a result, then say \emph{hētu} again
\end{itemize}

\begin{tcolorbox}[breakable,colback=teal!4,colframe=teal!35!black,title={Before you move on}]
What does \ml{ഹേതു} mean? (\textquotedblleft{}A cause\textquotedblright{}.) Say it once more.
\end{tcolorbox}
\section[\texorpdfstring{\ml{നിമിത്തം} (nimittaṁ)}{നിമിത്തം (nimittaṁ)}]{\ml{നിമിത്തം} (nimittaṁ) — on account of, because of}

\label{lesson:ML-C144-nimittam}



\begin{quote}
Before the new one: say the Malayalam for a result, then the Malayalam for a cause.
\end{quote}

\subsection*{You'll want to know: \ml{നിമിത്തം}}
\textbf{\ml{നിമിത്തം}} — \emph{nimittaṁ} — \textquotedblleft{}on account of, because of\textquotedblright{}.

\begin{grammarlens}[title={What you've built}]
One more word for this chapter.
\end{grammarlens}

\subsection*{Your turn}
\begin{itemize}
\raggedright
  \item \emph{Say it:} \emph{nimittaṁ}
  \item \emph{Say it:} \emph{nimittaṁ}, once more
  \item \emph{Say it:} say \emph{nimittaṁ}
  \item \emph{Recall:} say the Malayalam for a result, then the Malayalam for a cause, then say \emph{nimittaṁ} again
\end{itemize}

\begin{tcolorbox}[breakable,colback=teal!4,colframe=teal!35!black,title={Before you move on}]
What does \ml{നിമിത്തം} mean? (\textquotedblleft{}On account of, because of\textquotedblright{}.) Say it once more.
\end{tcolorbox}
